\subsection{Definition and first properties of extension modules}

For every A-module E, we denote by \(p_{\mathbf{E}} : \mathbf{L}(\mathbf{E}) \to \mathbf{E}\) (resp. \(e_{\mathbf{E}} : \mathbf{E} \to \mathbf{I}(\mathbf{E})\) ) the canonical free (resp. injective) resolution, cf. X, p. 50 (resp. p. 52).

DEFINITION 1. --- Let M and N be two A-modules. One calls the extension module of N by M the graded k-module

\[
\operatorname{Ext} _ {\mathbf {A}} (\mathbf {M}, \mathbf {N}) = \mathrm{H} \left(\operatorname{Homgr} _ {\mathbf {A}} (\mathrm{L} (\mathbf {M}), \mathrm{I} (\mathbf {N}))\right).\tag{4}
\]

The homogeneous components of Ext \(_{A}\) (M, N) are denoted

\[
\operatorname{Ext} _ {\mathbf {A}} ^ {n} (\mathbf {M}, \mathbf {N}) = \mathrm{H} ^ {n} \left(\operatorname{Homgr} _ {\mathbf {A}} (\mathbf {L} (\mathbf {M}), \mathrm{I} (\mathbf {N}))\right).\tag{5}
\]

Since L(M) (resp. I(N)) is zero on the right (resp. left), we have

\[
\operatorname{Ext} _ {\mathbf {A}} ^ {n} (\mathbf {M}, \mathbf {N}) = 0 \quad \text { for } \quad n <   0.\tag{6}
\]

Remark 1. --- We shall see below (X, p. 107, prop. 6) finiteness properties of the modules Ext \(_{A}\) (M, N). For example, if A is commutative Noetherian, and if M and N are finitely generated A-modules, each A-module Ext \(_{A}^{n}\) (M, N) is finitely generated.

Let \(f: \mathbf{M}' \to \mathbf{M}\) and \(g: \mathbf{N} \to \mathbf{N}'\) homomorphisms of A-modules, we set

\[
\operatorname{Ext} _ {\mathbf {A}} (f, g) = \mathrm{H} \left(\operatorname{Homgr} _ {\mathbf {A}} (\mathrm{L} (f), \mathrm{I} (g))\right);
\]

it is a degree 0 homomorphism of graded k-modules

\[
\operatorname{Ext} _ {\mathbf {A}} (f, g): \operatorname{Ext} _ {\mathbf {A}} (\mathbf {M}, \mathbf {N}) \rightarrow \operatorname{Ext} _ {\mathbf {A}} \left(\mathbf {M} ^ {\prime}, \mathbf {N} ^ {\prime}\right),
\]

whose homogeneous components are denoted

\[
\operatorname{Ext} _ {\mathbf {A}} ^ {n} (f, g): \operatorname{Ext} _ {\mathbf {A}} ^ {n} (\mathbf {M}, \mathbf {N}) \rightarrow \operatorname{Ext} _ {\mathbf {A}} ^ {n} \left(\mathbf {M} ^ {\prime}, \mathbf {N} ^ {\prime}\right).
\]

By prop. 1 of X, p. 82, the canonical homomorphism

\[
\lambda^ {0} (\mathrm{L} (\mathrm{M}), \mathrm{I} (\mathrm{N})): \mathrm{H} ^ {0} \left(\operatorname{Homgr} _ {\mathrm{A}} (\mathrm{L} (\mathrm{M}), \mathrm{I} (\mathrm{N}))\right)\rightarrow \operatorname{Hom} _ {\mathrm{A}} \left(\mathrm{H} _ {0} (\mathrm{L} (\mathrm{M})), \mathrm{H} ^ {0} (\mathrm{I} (\mathrm{N}))\right)
\]

is bijective; using the isomorphisms \(\mathbf{M} \to \mathrm{H}_{0}(\mathrm{L}(\mathbf{M}))\) and \(\mathrm{H}^{0}(\mathrm{I}(\mathbf{N})) \to \mathbf{N}\) , we deduce an isomorphism said to be canonical

\[
\lambda_ {\mathbf {M}, \mathbf {N}}: \operatorname{Ext} _ {\mathbf {A}} ^ {0} (\mathbf {M}, \mathbf {N}) \rightarrow \operatorname{Hom} _ {\mathbf {A}} (\mathbf {M}, \mathbf {N}).\tag{7}
\]

One will always identify \(\operatorname{Ext}_{\mathbf{A}}^{0}(\mathbf{M},\mathbf{N})\) with \(\operatorname{Hom}_{\mathbf{A}}(\mathbf{M},\mathbf{N})\) by this isomorphism. Then the map \(k\) -linear \(\operatorname{Ext}_{\mathbf{A}}^{0}(f,g)\) is identified with \(\operatorname{Hom}_{\mathbf{A}}(f,g)\) .

Remark 2. --- The morphism of complexes

\[
\operatorname{Homgr} _ {\mathbf {A}} \left(p _ {\mathbf {M}}, e _ {\mathbf {N}}\right): \operatorname{Hom} _ {\mathbf {A}} (\mathbf {M}, \mathbf {N}) \rightarrow \operatorname{Homgr} _ {\mathbf {A}} \left(\mathbf {L} (\mathbf {M}), \mathrm{I} (\mathbf {N})\right)
\]

induces on the homology of degree 0 the isomorphism

\[
\lambda_ {\mathbf {M}, \mathbf {N}} ^ {- 1}: \operatorname{Hom} _ {\mathbf {A}} (\mathbf {M}, \mathbf {N}) \rightarrow \operatorname{Ext} _ {\mathbf {A}} ^ {0} (\mathbf {M}, \mathbf {N})
\]

inverse of \(\lambda_{M,N}\) .

One has \(L(1_{\mathbf{M}}) = 1_{L(\mathbf{M})}\) , \(I(1_{\mathbf{N}}) = 1_{I(\mathbf{N})}\) , hence by passing to homology

\[
\operatorname{Ext} _ {\mathbf {A}} \left(1 _ {\mathbf {M}}, 1 _ {\mathbf {N}}\right) = 1 _ {\operatorname{Ext} _ {\mathbf {A}} (\mathbf {M}, \mathbf {N})}.\tag{8}
\]

If \(f': M'' \to M'\) and \(g': N' \to N''\) are homomorphisms of A-modules, we have \(L(f \circ f') = L(f) \circ L(f')\) and \(I(g' \circ g) = I(g') \circ I(g)\) , hence

\[
\operatorname{Ext} _ {\mathbf {A}} \left(f \circ f ^ {\prime}, g ^ {\prime} \circ g\right) = \operatorname{Ext} _ {\mathbf {A}} \left(f ^ {\prime}, g ^ {\prime}\right) \circ \operatorname{Ext} _ {\mathbf {A}} (f, g).\tag{9}
\]

Consider the morphisms of k-complexes

\[
\operatorname{Homgr} _ {\mathbf {A}} (\mathrm{L} (\mathrm{M}), \mathrm{N}) \xrightarrow {\operatorname{Homgr} _ {\mathbf {A}} (1 , e _ {\mathrm{N}})} \operatorname{Homgr} _ {\mathbf {A}} (\mathrm{L} (\mathrm{M}), \mathrm{I} (\mathrm{N})) \xleftarrow {\operatorname{Homgr} _ {\mathbf {A}} (p _ {\mathrm{M}}, 1)} \operatorname{Homgr} _ {\mathbf {A}} (\mathrm{M}, \mathrm{I} (\mathrm{N})),
\]

and the homomorphisms they induce in homology

\[
\mathrm{H} \left(\operatorname{Homgr} _ {\mathrm{A}} (\mathrm{L} (\mathrm{M}), \mathrm{N})\right) \xrightarrow {\varphi_ {\mathrm{M}} (\mathrm{N})} \operatorname{Ext} _ {\mathrm{A}} (\mathrm{M}, \mathrm{N}) \xleftarrow {\overline {{\varphi}} _ {\mathrm{N}} (\mathrm{M})} \mathrm{H} \left(\operatorname{Homgr} _ {\mathrm{A}} (\mathrm{M}, \mathrm{I} (\mathrm{N}))\right).
\]

By prop. 4 of X, p. 86, \(\operatorname{Homgr}_{\mathbf{A}}(1, e_{\mathbf{N}})\) and \(\operatorname{Homgr}_{\mathbf{A}}(p_{\mathbf{M}}, 1)\) are homologisms, hence:

PROPOSITION 5. --- The k-homomorphisms

\[
\begin{array}{l l} & \varphi_ {\mathbf {M}} (\mathbf {N}): \mathrm{H} (\mathrm{Homgr} _ {\mathbf {A}}   (\mathrm{L} (\mathbf {M}), \mathrm{N})) \to \mathrm{Ext} _ {\mathbf {A}}   (\mathbf {M}, \mathbf {N}) \\ 	ext{and} & \overline {{\varphi}} _ {\mathbf {N}} (\mathbf {M}): \mathrm{H} (\mathrm{Homgr} _ {\mathbf {A}}   (\mathbf {M}, \mathrm{I} (\mathbf {N}))) \to \mathrm{Ext} _ {\mathbf {A}}   (\mathbf {M}, \mathbf {N}) 	ext{are bijective.} \end{array}
\]

COROLLARY. --- If M is projective (resp. if N is injective), then Ext \(_{A}^{i}\) (M, N) = 0 for i \textgreater{} 0.

Indeed,

\[
\operatorname{Homgr} _ {\mathbf {A}} \left(1, e _ {\mathbf {N}}\right): \operatorname{Hom} _ {\mathbf {A}} (\mathbf {M}, \mathbf {N}) \rightarrow \operatorname{Homgr} _ {\mathbf {A}} (\mathbf {M}, \mathrm{I} (\mathbf {N}))
\]

(resp. Homgr \(_{A}\) ( \(p_{M}\) , 1): Hom \(_{A}\) (M, N) → Homgr \(_{A}\) (L(M), N)) is then a homologism (X, p. 86, prop. 4), whence the conclusion.

Remarks. --- 3) If \(g: \mathbf{N} \to \mathbf{N}'\) is a homomorphism of A-modules, then

\[
\operatorname{Homgr} _ {\mathbf {A}} \left(1 _ {\mathrm{L} (\mathbf {M})}, g\right) \circ \operatorname{Homgr} _ {\mathbf {A}} \left(1 _ {\mathrm{L} (\mathbf {M})}, e _ {\mathbf {N}}\right) = \operatorname{Homgr} _ {\mathbf {A}} \left(1 _ {\mathrm{L} (\mathbf {M})}, e _ {\mathbf {N} ^ {\prime}}\right) \circ \operatorname{Homgr} _ {\mathbf {A}} \left(1 _ {\mathrm{L} (\mathbf {M})}, \mathrm{I} (g)\right)
\]

therefore the diagram

\[
\begin{array}{c} \mathrm{H} (\operatorname{Homgr} _ {\mathbf {A}} (\mathrm{L} (\mathbf {M}), \mathrm{N})) \xrightarrow {\varphi_ {\mathbf {M}} (\mathrm{N})} \operatorname{Ext} _ {\mathbf {A}} (\mathbf {M}, \mathbf {N}) \\ \mathrm{H} (\operatorname{Homgr} _ {\mathbf {A}} (1 _ {\mathrm{L} (\mathbf {M})}, g)) \Bigg | _ {\downarrow} \qquad \qquad \qquad \qquad \qquad \qquad \qquad \qquad \qquad \qquad \qquad \qquad \qquad \qquad \qquad \qquad \qquad \qquad \qquad \qquad \qquad \qquad \qquad \qquad \qquad \qquad \qquad \qquad \qquad \qquad \qquad \qquad \qquad \qquad \\ \mathrm{H} (\operatorname{Homgr} _ {\mathbf {A}} (\mathrm{L} (\mathbf {M}), \mathrm{N} ^ {\prime})) \xrightarrow {\varphi_ {\mathbf {M}} (\mathrm{N} ^ {\prime})} \operatorname{Ext} _ {\mathbf {A}} (\mathbf {M}, \mathrm{N} ^ {\prime}) \end{array}
\]

is commutative.

\begin{enumerate}
\def\labelenumi{\arabic{enumi})}
\setcounter{enumi}{3}
\tightlist
\item
  Likewise, if \(f: \mathbf{M}' \to \mathbf{M}\) is a homomorphism of A-modules, the diagram
\end{enumerate}

\[
\begin{array}{c c c} \mathrm{H} (\mathrm{Homgr} _ {\mathrm{A}} (\mathrm{M}, \mathrm{I} (\mathrm{N}))) & \xrightarrow {\overline {{\varphi}} _ {\mathrm{N}} (\mathrm{M})} & \mathrm{Ext} _ {\mathrm{A}} (\mathrm{M}, \mathrm{N}) \\ \mathrm{H} (\mathrm{Homgr} _ {\mathrm{A}} (f, 1 _ {\mathrm{I} (\mathrm{N})})) \Big \downarrow & & \Big \downarrow \mathrm{Ext} _ {\mathrm{A}} (f, 1 _ {\mathrm{N}}) \\ \mathrm{H} (\mathrm{Homgr} _ {\mathrm{A}} (\mathrm{M} ^ {\prime}, \mathrm{I} (\mathrm{N}))) & \xrightarrow {\overline {{\varphi}} _ {\mathrm{N}} (\mathrm{M} ^ {\prime})} & \mathrm{Ext} _ {\mathrm{A}} (\mathrm{M} ^ {\prime}, \mathrm{N}) \end{array}
\]

is commutative.

PROPOSITION 6. --- The map (f, g) ↦ Ext\_A (f, g) :

\[
\operatorname{Hom} _ {\mathbf {A}} \left(\mathrm{M} ^ {\prime}, \mathrm{M}\right) \times \operatorname{Hom} _ {\mathbf {A}} \left(\mathrm{N}, \mathrm{N} ^ {\prime}\right)\rightarrow \operatorname{Homgr} _ {k} \left(\operatorname{Ext} _ {\mathbf {A}} (\mathrm{M}, \mathrm{N}), \operatorname{Ext} _ {\mathbf {A}} \left(\mathrm{M} ^ {\prime}, \mathrm{N} ^ {\prime}\right)\right)
\]

is k-bilinear.

Let \(f \in \operatorname{Hom}_{\mathbf{A}}(\mathbf{M}', \mathbf{M})\) , \(g_1, g_2 \in \operatorname{Hom}_{\mathbf{A}}(\mathbf{N}, \mathbf{N}')\) , \(\lambda_1, \lambda_2 \in k\) ; then the morphisms \(\operatorname{Homgr}_{\mathbf{A}}(\mathbf{L}(f), \lambda_1 g_1 + \lambda_2 g_2)\) and \(\lambda_1 \operatorname{Homgr}_{\mathbf{A}}(\mathbf{L}(f), g_1) + \lambda_2 \operatorname{Homgr}_{\mathbf{A}}(\mathbf{L}(f), g_2)\) of \(\operatorname{Homgr}_{\mathbf{A}}(\mathbf{L}(\mathbf{M}), \mathbf{N})\) to \(\operatorname{Homgr}_{\mathbf{A}}(\mathbf{L}(\mathbf{M}), \mathbf{N}')\) coincide; therefore, by prop. 5 and remark 3,

\[
\operatorname{Ext} _ {\mathbf {A}} \left(f, \lambda_ {1} g _ {1} + \lambda_ {2} g _ {2}\right) = \lambda_ {1} \operatorname{Ext} _ {\mathbf {A}} \left(f, g _ {1}\right) + \lambda_ {2} \operatorname{Ext} _ {\mathbf {A}} \left(f, g _ {2}\right).\tag{10}
\]

We reason similarly for the mapping. \(f\mapsto\mathrm{Ext}_{\mathbf{A}}(f,g)\) .

COROLLARY. --- Let \(\lambda \in k\). If \(\lambda\) annihilates M or N, it annihilates Ext \(_{A}\) (M, N).

Indeed, \(\lambda 1_{\mathrm{Ext}(\mathbf{M},\mathbf{N})} = \mathrm{Ext}(\lambda 1_{\mathbf{M}},1_{\mathbf{N}}) = \mathrm{Ext}(1_{\mathbf{M}},\lambda 1_{\mathbf{N}})\) .

PROPOSITION 7. --- Let I and J be two sets, \((\mathbf{M}_{\alpha})_{\alpha \in I}\) and \((\mathbf{N}_{\beta})_{\beta \in J}\) two families of A-modules; the homomorphism \(\operatorname{Ext}_{\mathbf{A}}\left(\bigoplus_{\alpha \in I} \mathbf{M}_{\alpha}, \prod_{\beta \in J} \mathbf{N}_{\beta}\right) \to \prod_{\beta \in J, \alpha \in I} \operatorname{Ext}_{\mathbf{A}}(\mathbf{M}_{\alpha}, \mathbf{N}_{\beta})\) deduced from the canonical homomorphisms \(\mathbf{M}_{\alpha} \to \oplus \mathbf{M}_{\alpha}\) and \(\Pi \mathbf{N}_{\beta} \to \mathbf{N}_{\beta}\) is bijective.

It suffices to prove that for every A-module M (resp. N), the homomorphisms

\[
\operatorname{Ext} \left(\mathrm{M}, \Pi \mathrm{N} _ {\beta}\right)\rightarrow \Pi \operatorname{Ext} \left(\mathrm{M}, \mathrm{N} _ {\beta}\right) \quad (\text { resp.   Ext } (\oplus \mathrm{M} _ {\alpha}, \mathrm{N}) \rightarrow \Pi \operatorname{Ext} \left(\mathrm{M} _ {\alpha}, \mathrm{N}\right))
\]

are bijective. Now this follows from what precedes, from prop. 1 of X, p. 28, and from the canonical isomorphisms \(\operatorname{Homgr}_{\mathbf{A}}(\mathbf{L}(\mathbf{M}), \Pi \mathbf{N}_{\beta}) \to \Pi \operatorname{Homgr}_{\mathbf{A}}(\mathbf{L}(\mathbf{M}), \mathbf{N}_{\beta})\) and \(\operatorname{Homgr}_{\mathbf{A}}(\oplus \mathbf{M}_{\alpha}, \mathrm{I}(\mathbf{N})) \to \Pi \operatorname{Homgr}_{\mathbf{A}}(\mathbf{M}_{\alpha}, \mathrm{I}(\mathbf{N}))\) .

Remark 5. --- Let P and Q be two right A-modules. One defines Ext \(_{A}\) (P, Q) by

\[
\begin{array}{r l} \operatorname{Ext} _ {\mathbf {A}} (\mathrm{P}, \mathrm{Q}) & = \mathrm{H} (\operatorname{Homgr} _ {\mathbf {A}} (\mathrm{L} (\mathrm{P}), \mathrm{I} (\mathrm{Q}))) = \mathrm{H} (\operatorname{Homgr} _ {\mathbf {A} ^ {\circ}} (\mathrm{L} (\mathrm{P} ^ {\circ}), \mathrm{I} (\mathrm{Q} ^ {\circ}))) = \\ & = \operatorname{Ext} _ {\mathbf {A} ^ {\circ}} (\mathrm{P} ^ {\circ}, \mathrm{Q} ^ {\circ}). \end{array}
\]

All definitions and propositions in this paragraph therefore apply to right A-modules by viewing them as left modules over the ring A°.

